%% NeurIPS-style submission, self-contained preamble (no external .sty required).
%% Compile with: pdflatex Project_Proposal.tex (twice for cross-references).
\documentclass[11pt]{article}

%% ---------- Geometry & typography (NeurIPS-like) ----------
\usepackage[letterpaper, top=1.0in, bottom=1.0in, left=1.25in, right=1.25in]{geometry}
\usepackage[T1]{fontenc}
\usepackage[utf8]{inputenc}
\usepackage{times}
\usepackage{microtype}
\usepackage{setspace}\setstretch{1.02}
\usepackage{parskip}\setlength{\parskip}{3pt}

%% ---------- Math, graphics, tables ----------
\usepackage{amsmath, amssymb, amsthm, mathtools}
\usepackage{graphicx}
\usepackage{booktabs}
\usepackage{multirow}
\usepackage{array}
\usepackage{enumitem}
\usepackage{caption}\captionsetup{font=small, labelfont=bf, skip=4pt}
\usepackage{subcaption}
\usepackage{float}

%% ---------- Hyperlinks & cross-refs ----------
\usepackage{hyperref}
\hypersetup{colorlinks=true, linkcolor=black, citecolor=black, urlcolor=black,
            pdftitle={OPTFM is 1-WL bounded on bipartite MILPs},
            pdfauthor={Md Abrar Jahin and Shahriar Tanvir Alam}}
\usepackage[capitalise,noabbrev]{cleveref}

%% ---------- Sections (NeurIPS-style flush left, sans-bold) ----------
\usepackage{titlesec}
\titleformat{\section}{\large\bfseries}{\thesection}{0.6em}{}
\titlespacing*{\section}{0pt}{12pt}{4pt}
\titleformat{\subsection}{\normalsize\bfseries}{\thesubsection}{0.5em}{}
\titlespacing*{\subsection}{0pt}{8pt}{2pt}
\titleformat{\subsubsection}[runin]{\itshape}{\thesubsubsection}{0.4em}{}
\titlespacing*{\subsubsection}{0pt}{4pt}{0.6em}

%% ---------- Theorem environments ----------
\theoremstyle{plain}
\newtheorem{theorem}{Theorem}
\newtheorem{lemma}{Lemma}
\newtheorem{proposition}{Proposition}
\newtheorem{corollary}{Corollary}
\theoremstyle{definition}
\newtheorem{definition}{Definition}
\theoremstyle{remark}
\newtheorem*{remark}{Remark}

%% ---------- Listings (used only briefly in appendix) ----------
\usepackage{listings}
\lstset{basicstyle=\ttfamily\footnotesize, breaklines=true,
        frame=single, framerule=0.3pt, xleftmargin=4pt, columns=fullflexible}

%% ---------- Macros ----------
\newcommand{\optfm}{\textsc{OPTFM}}
\newcommand{\sgformer}{\textsc{SGFormer}}
\newcommand{\rwpe}{\textsc{RWPE}}
\newcommand{\rni}{\textsc{RNI}}
\newcommand{\given}{\,\vert\,}
\newcommand{\multiset}[1]{\{\!\!\{#1\}\!\!\}}
\DeclareMathOperator*{\Pool}{Pool}
\DeclareMathOperator*{\HASH}{HASH}
\DeclareMathOperator*{\LinAttn}{LinAttn}

\title{\vspace{-1.3cm}\textbf{\Large OPTFM is 1-WL Bounded on Bipartite MILPs:\\
A Theorem, a Two-Family Empirical Probe, and a Structural Probe Battery}}

\author{
\textbf{Md Abrar Jahin}\;\,\textbf{Shahriar Tanvir Alam}\\
University of Southern California\\
\texttt{\{mdabrar, shahriar\}@usc.edu}
}

\date{}

\begin{document}
\maketitle
\thispagestyle{empty}

\begin{abstract}
\noindent
We investigate whether \optfm{}~\cite{yuan2025optfm}---the first graph foundation
model for combinatorial optimization---escapes the 1-Weisfeiler--Lehman (1-WL)
expressiveness bound of Chen et al.~\cite{chen2023expressive} for bipartite
mixed-integer linear programs (MILPs). \optfm{} motivates its hierarchical
multi-view cross-attention as ``going beyond the local interaction structure of
standard GNNs.'' We show this is not the case. \textbf{Theorem.} The full
hierarchical \optfm{} encoder (the \texttt{SGFormer\_MIP\_hierarchical\_improve}
class shipped in \optfm{}'s source) is 1-WL bounded for \emph{any} weight
setting; the bound is preserved by \optfm{}'s virtual-global-node pretraining
augmentation. \textbf{Empirical probe.} On 15 verified 1-WL-equivalent
non-isomorphic bipartite MILP pairs from two structurally independent families
(2-regular cycles \(C_{4k}\) vs \(k\cdot C_4\) for \(k\!\in\!\{2,\dots,30\}\),
plus a 3-regular cubic pair), six architectures including the random-init full
hierarchical \optfm{} produce \emph{bit-identical} graph embeddings on every
pair. Layer-wise instrumentation confirms the bound holds at every sub-layer.
A 2-class supervised fine-tune of backbone+head pins cross-entropy at
\(\ln 2\!\approx\!0.6931\) for 300 epochs across seeds---training cannot
escape an architectural bound. \textbf{Probe battery.} An MLP probe on frozen
pretrained embeddings collapses from \(R^2\!=\!{+}0.387\) on random MILPs to
\(R^2\!=\!{-}0.011\) on 1-WL-equivalent MILPs for \texttt{n\_components},
\emph{despite the target having six times larger variance} on the latter.
A minimal Random Walk Positional Encoding transform breaks the bound across
all architectures, confirming the probe is genuinely falsifiable.
\end{abstract}

\section{Introduction}
\label{sec:intro}

Graph neural networks (GNNs) for mixed-integer programming have been studied
extensively as policy or representation networks for branch-and-bound and
its surrogates~\citep{gasse2019exact, khalil2017learning, gupta2020hybrid}.
Chen et al.~\cite{chen2023expressive} established that any message-passing
GNN applied to a bipartite-graph encoding of a MILP is at most as
discriminative as 1-Weisfeiler--Lehman (1-WL) color refinement. They
exhibited a pair of feature-labeled bipartite MILPs that are 1-WL equivalent
but differ in feasibility---demonstrating that the bound is tight for the
most basic decision in optimization.

\optfm{}~\cite{yuan2025optfm}, the first foundation model for
combinatorial optimization, proposes a hierarchical multi-view graph
transformer specifically motivated as going \emph{beyond} GNN-style local
interaction via global linear attention and edge-weighted cross-attention
between constraint and variable nodes. The OPTFM paper claims that this
architecture ``enhances representational power,'' but provides neither an
expressiveness theorem nor a test on known indistinguishable pairs.

\paragraph{Question.} Does \optfm{}'s hierarchical encoder distinguish 1-WL
equivalent non-isomorphic bipartite MILPs?

\paragraph{Contributions.}
\begin{enumerate}[leftmargin=18pt, topsep=0pt, itemsep=1pt]
    \item \textbf{Theorem (\Cref{sec:theorem}).} For any weights, \optfm{}'s
    hierarchical encoder factors through 1-WL equivalence. The proof composes
    five lemmas covering \sgformer{}-style linear attention~\cite{wu2023sgformer},
    \optfm{}'s additive edge-weighted cross-attention, the bipartite GCN
    branch~\cite{gasse2019exact}, mean/sum/max pooling, and the
    virtual-global-node pretraining augmentation.

    \item \textbf{Two-family empirical probe (\Cref{sec:main}).} On 15 verified
    1-WL-equivalent non-isomorphic bipartite MILP pairs drawn from two
    structurally independent families (2-regular cycles and a 3-regular cubic
    pair), six architectures including the full hierarchical \optfm{} produce
    bit-identical pooled embeddings on every pair. Layer-wise instrumentation
    confirms the bound stage-by-stage (\Cref{fig:layerwise}). A supervised
    fine-tune of the backbone plus a 2-class linear head pins the cross-entropy
    loss at $\ln 2 \approx 0.6931$ for 300 epochs across three seeds
    (\Cref{fig:plateau})---training cannot escape an architectural bound.

    \item \textbf{Probe battery (\Cref{sec:probes}).} On 800 random bipartite
    MILPs, an MLP probe on frozen pretrained \optfm{} embeddings predicts
    \texttt{n\_components} at $R^2={+}0.387$. On 420 1-WL-equivalent MILPs the
    same probe collapses to $R^2={-}0.011$, despite a $6{\times}$ larger target
    variance. The \texttt{lp\_value} target is at chance on both populations,
    indicating that the masked-edge pretraining objective does not induce
    LP-aware representations.

    \item \textbf{Falsifiability (\Cref{sec:rwpe}).} Injecting Random Walk
    Positional Encodings~\cite{dwivedi2022rwpe} into three previously-zero
    Ecole feature slots breaks the bit-identity for every architecture and
    family, ruling out the ``pathological pair'' defense.
\end{enumerate}

In the course of this investigation we also discovered that the ``canonical
1-WL pair'' of variable-feature-equivariant bipartite MILPs widely used as a
sanity-check in earlier probing code is in fact \emph{graph-isomorphic} as a
feature-labeled bipartite graph; we document this artifact in
\Cref{app:broken-canonical} and built our probe families to be free of it.

\section{Preliminaries}
\label{sec:prelim}

\paragraph{Feature-labeled bipartite MILP.}
Following~\cite{gasse2019exact, chen2023expressive}, a feature-labeled
bipartite MILP is a tuple $M = (C, V, A, f_C, f_V)$ with constraint nodes $C$,
variable nodes $V$, coefficient matrix $A \in \mathbb R^{|C|\times|V|}$, and
per-node features $f_C, f_V$. We use Ecole's \texttt{MilpBipartite}
extractor: $d_v=9$ variable features (objective coefficient, integrality
flag, bound flags and values, basis status, reduced cost, solution value)
and $d_c=1$ constraint feature (normalized RHS $b_i / \|A_i\|_2$). The
\optfm{} pretrained checkpoint was trained on these exact features.

\paragraph{1-WL refinement.}
For a feature-labeled bipartite MILP $M$, the iterative coloring is
\begin{equation}\label{eq:wl}
\chi^{(0)}(u) = f(u),\quad \chi^{(t+1)}(u) = \HASH\!\Big(\chi^{(t)}(u),\,
   \multiset{(\chi^{(t)}(w), A_{uw}) : w \sim u}\Big),
\end{equation}
where $\multiset{\cdot}$ denotes a multiset. Two MILPs $M_1, M_2$ are
\emph{1-WL equivalent}, written $M_1 \equiv_{1\text{WL}} M_2$, if for every
$t$ the global color multisets coincide. Isomorphism implies 1-WL
equivalence; the converse is false in general~\cite{cai1992lower}.

\paragraph{The hierarchical \optfm{} encoder.}
The class \texttt{SGFormer\_MIP\_hierarchical\_improve}
in \texttt{OPTFM/node\_pretrain/ours\_crossattention\_improve.py} maps a
bipartite MILP to a graph embedding $z\in\mathbb R^h$ via
(i) input embeddings $\tilde V_0 = W_v f_V$, $\tilde C_0 = W_c f_C$;
(ii) intra-type \sgformer{} self-attention on each side independently;
(iii) cross-attention with edge-weighted aggregation, defined for the $C\to V$
pass as
\begin{equation}\label{eq:cross}
\tilde V_2 = W_o \big[\,\LinAttn(\tilde V_1, \tilde C_1);\;
                          A^\top \tilde C_1\,\big],
\end{equation}
and symmetrically $V\to C$;
(iv) a parallel bipartite GCN branch
$(\tilde V^{(g)},\tilde C^{(g)}) = \texttt{GNNPolicy}(\cdots)$
following~\cite{gasse2019exact};
(v) convex fuse $\tilde V = \alpha \tilde V^{(g)} + (1-\alpha) \tilde V_2$
and analogously for $\tilde C$;
(vi) pooling $z = \Pool(\tilde C \oplus \tilde V)$ with
$\Pool\in\{\text{mean},\text{sum},\text{max}\}$.

The \sgformer{} linear-attention block~\cite{wu2023sgformer} computes for each
query index $i$
\begin{equation}\label{eq:linattn}
\LinAttn(Q,X)_i = D_i^{-1}\!\Big(V_i + \tfrac{1}{N}\,\tilde Q_i\,\big(\tilde K^\top V\big)\Big),
\quad D_i = 1 + \tfrac{1}{N}\tilde Q_i (\tilde K^\top \mathbf 1),
\end{equation}
with Frobenius-norm normalized $\tilde Q = W_q Q / \|W_q Q\|_F$ and
$\tilde K = K/\|K\|_F$. The pretraining pipeline in
\texttt{OPTFM/node\_pretrain/main\_mip.py}, lines 148--171, additionally
appends two virtual global nodes per subgraph: a constraint-side global with
feature equal to the mean of constraint features (connected to every variable
with unit weight) and a variable-side analogue.

\section{OPTFM is 1-WL bounded}
\label{sec:theorem}

\begin{theorem}[\optfm{} hierarchical encoder is 1-WL bounded]
\label{thm:main}
Let $M_1, M_2$ be feature-labeled bipartite MILPs with $M_1 \equiv_{1\text{WL}} M_2$.
For \emph{any} weights $\theta$ of the hierarchical \optfm{} encoder
\(\Phi_\theta = \texttt{SGFormer\_MIP\_hierarchical\_improve}\) and any pooling
in $\{\text{mean},\text{sum},\text{max}\}$,
\begin{equation*}
\Phi_\theta(M_1) \;=\; \Phi_\theta(M_2).
\end{equation*}
The same conclusion holds when both inputs are augmented with the
virtual-global-node construction from \optfm{}'s pretraining pipeline.
\end{theorem}

\begin{corollary}\label{cor:concrete}
\optfm{} cannot distinguish $C_{4k}$-bipartite from $k\cdot C_4$-bipartite for
any $k\!\ge\!2$, nor the connected $6{+}6$ cubic bipartite graph from
$2\cdot K_{3,3}$, despite both pairs being non-isomorphic.
\end{corollary}

\subsection{Lemmas}

We say an operator $\Phi$ is \emph{1-WL bounded} if multiset-equal node
inputs produce multiset-equal node outputs.

\begin{lemma}[Linear attention is a multiset function of the source]\label{lem:linattn}
For \(Q,K,V\in\mathbb R^{N\times d}\) arising from a node-wise linear map,
\(\LinAttn(Q,X)_i\) defined in~\Cref{eq:linattn} depends only on $(q_i, v_i)$
and on the multisets $\multiset{k_\ell}$ and $\multiset{(k_\ell, v_\ell)}$.
\end{lemma}
\begin{proof}
Both $\sum_\ell k_\ell^\top v_\ell$ and $\sum_\ell k_\ell^\top \mathbf 1$ are
commutative sums over $\ell$, hence symmetric functions of the multiset of
source pairs. The Frobenius normalizer $\|Q\|_F^2 = \sum_j \|q_j\|_2^2$ is a
sum of per-node quantities. LayerNorm, residual, ReLU, and dropout (identity
at eval) are per-node. \qedhere
\end{proof}

\begin{lemma}[Cross-attention with edge-weighted aggregation]\label{lem:cross}
The block \texttt{TransConvCross} that linearly attends across types and
concatenates the dense edge-weighted message $A^\top \tilde c$ before a linear
projection (\Cref{eq:cross}) is 1-WL bounded.
\end{lemma}
\begin{proof}
Linear attention is bounded by~\Cref{lem:linattn}. The aggregation
$m_j = \sum_i A_{ij}\,\tilde c_i$ is a sum over the edge-weighted neighborhood
$\multiset{(A_{ij},\tilde c_i)}$, which 1-WL equivalence preserves by
construction. Concatenation and a per-node linear projection commute with
multiset equivalence. The variable-to-constraint pass is symmetric. \qedhere
\end{proof}

\begin{lemma}[The bipartite GCN branch]\label{lem:gcn}
The branch \texttt{GNNPolicy} in both the baseline and hierarchical
\optfm{} variants is 1-WL bounded on graph-level outputs.
\end{lemma}
\begin{proof}
\texttt{GNNPolicy} is a Gasse-style bipartite MPNN~\cite{gasse2019exact} with
two convolutional layers; it instantiates the MPNN class for which Xu et
al.~\cite{xu2019powerful} and Morris et al.~\cite{morris2019wlgo} proved the
1-WL bound, and which Chen et al.~\cite{chen2023expressive} specialized to
bipartite MILPs. \qedhere
\end{proof}

\begin{lemma}[Convex fuse and pooling]\label{lem:fuse}
The fuse $\tilde v \mapsto \alpha\tilde v^{(g)}+(1-\alpha)\tilde v$ (and
analogously for $\tilde c$) is per-node, and mean/sum/max pooling are
symmetric functions of their input multiset; both preserve 1-WL bounds.
\end{lemma}

\begin{lemma}[Virtual global nodes preserve 1-WL]\label{lem:virtnode}
Let $T$ append a global constraint node with feature
$\bar f_C = |C|^{-1}\sum_c f_C(c)$, a global variable node with feature
$\bar f_V = |V|^{-1}\sum_v f_V(v)$, unit-weight edges from the global
constraint to every variable, and unit-weight edges from the global variable
to every constraint. If $M_1 \equiv_{1\text{WL}} M_2$, then
$T(M_1) \equiv_{1\text{WL}} T(M_2)$.
\end{lemma}
\begin{proof}
At iteration 0, $\bar f_C(M_1)=\bar f_C(M_2)$ and similarly for $\bar f_V$
because the original feature multisets are equal by hypothesis. At iteration
$t{+}1$, every original node's color refinement~(\Cref{eq:wl}) gains the
single neighborhood term $(\chi^{(t)}(\text{global}), 1)$, which is identical
in $M_1$ and $M_2$ by induction on $t$. The global nodes themselves see the
entire (multiset-equal) opposite-side neighborhood and therefore share the
same color across the two graphs. \qedhere
\end{proof}

\subsection{Proof of \Cref{thm:main}}

Apply~\Cref{lem:virtnode} to optionally augment with virtual global nodes;
1-WL equivalence is preserved. Propagate stage-by-stage through the encoder:
input embeddings (per-node, trivially preserve multiset equivalence);
\texttt{TransConvSelf} (\Cref{lem:linattn}); \texttt{TransConvCross}
(\Cref{lem:cross}); \texttt{GNNPolicy} (\Cref{lem:gcn}); convex fuse and
pooling (\Cref{lem:fuse}). At every stage, multiset equivalence is the
loop invariant. Hence $\Phi_\theta(M_1) = \Phi_\theta(M_2)$. \qed

\paragraph{Remark on scope.}
The theorem is a strong, quantifier-over-weights statement: no
optimization trajectory through $\Theta$ can find a weight setting that
distinguishes a 1-WL equivalent pair. The bound is escapable only by input
transformations that are not themselves 1-WL invariants---random node
features~\cite{abboud2020rni, sato2021random} or positional
encodings~\cite{dwivedi2022rwpe}. We exploit this in~\Cref{sec:rwpe} as the
positive control of our probe.

\section{Empirical setup}
\label{sec:setup}

\paragraph{Pair families.}
We use two structurally independent families of feature-labeled bipartite
MILP pairs verified to be 1-WL equivalent and non-isomorphic.
\textbf{Family I (2-regular cycles)}, $k\!\in\!\{2,3,4,5,6,7,8,9,10,12,15,20,25,30\}$:
$G_A(k) = C_{4k}$-bipartite is a single $4k$-cycle alternating constraint and
variable nodes; $G_B(k) = k\cdot C_4$-bipartite is $k$ disjoint
$K_{2,2}$'s. Both have $2k$ constraints, $2k$ variables, all of degree 2.
With uniform $b\!=\!1$ the Ecole RHS feature is $1/\sqrt 2$ uniformly,
forcing 1-WL refinement to a single color class. They differ only in
connectivity: $G_A$ has 1 component, $G_B(k)$ has $k$.
\textbf{Family II (3-regular cubic)}: $G_A$ is a connected 3-regular
bipartite graph on $6{+}6$ vertices built from the circulant offsets
$\{0,1,3\}\bmod 6$; $G_B = K_{3,3}\sqcup K_{3,3}$. 1-WL equivalence by
uniform cubic regularity; 1 vs 2 components.

Equivalence and non-isomorphism are verified by brute-force permutation
search for $|C|+|V|\!\le\!12$ and NetworkX VF2 feature-labeled bipartite
isomorphism checks for larger pairs (\texttt{data/milp\_pairs\_v2.diagnose\_pair}).
\Cref{fig:families} shows the two families.

\begin{figure}[t]
\centering
\includegraphics[width=0.92\linewidth]{results/figures/fig1_pair_families.pdf}
\caption{\textbf{Two structurally independent 1-WL-equivalent non-isomorphic
bipartite MILP pair families.} Constraints (circles) and variables (squares);
edges from the nonzero pattern of $A$. Nodes are colored by their 1-WL
fixed-point class---uniform features collapse every constraint into one class
(blue) and every variable into one class (orange) in both members of each
pair. (a) $G^{\mathrm I}_A$: a single $4k$-cycle ($k{=}3$ shown, 12 vertices,
connected). (b) $G^{\mathrm I}_B$: $k$ disjoint $C_4$'s ($k$ components).
(c) $G^{\mathrm{II}}_A$: connected 3-regular bipartite on $6{+}6$ vertices
from circulant offsets $\{0,1,3\}\bmod 6$. (d) $G^{\mathrm{II}}_B$: $K_{3,3}
\sqcup K_{3,3}$. Within each family, members share the 1-WL color histogram
at every depth but differ in number of connected components---a 1-WL-invisible
graph invariant.}
\label{fig:families}
\end{figure}

\paragraph{Models.}
Six architectures are probed: the pretrained \sgformer{}+GCN baseline
(\texttt{Models\_SCIP/model\_params\_epoch\_31.pth}, 6{,}820 parameters,
verified loaded with $\texttt{strict=True}$ in
\texttt{scripts/verify\_weights.py}); the same architecture at random
initialization; \texttt{TransConvOnly} (1{,}344 params); \texttt{GNNOnly}
(5{,}602 params); a 2-layer \texttt{SimpleGCN} (736 params); and the full
hierarchical \optfm{} (\texttt{HierarchicalOPTFM} in
\texttt{models/optfm\_hierarchical.py}, 10{,}852 params)---a dense-operation
port of \texttt{SGFormer\_MIP\_hierarchical\_improve} that does not require
\texttt{torch\_sparse} (which has no prebuilt wheel for recent PyTorch on
Windows). Since 1-WL boundedness is an architectural property
(\Cref{thm:main} quantifies over $\theta$), random initialization suffices to
probe the hierarchical variant.

\paragraph{Input transforms.}
Eight transforms registered in \texttt{scripts/improvements.TRANSFORMS}:
\texttt{baseline} (identity); \texttt{rnf\_sigma\_0.3}, \texttt{rnf\_sigma\_1.0}
(Gaussian random node features into the three zero Ecole variable slots,
$\sigma\!\in\!\{0.3, 1.0\}$, 30 i.i.d.\ samples per pair);
\texttt{lp\_primal}, \texttt{lp\_reduced}, \texttt{lp\_primal\_dual} (LP-derived
features); \texttt{rwpe\_steps\_4\_6\_8} (Random Walk Positional Encoding
return probabilities $P^L[i,i]$ at $L\!\in\!\{4,6,8\}$); and
\texttt{virtual\_global\_node} (the pretraining augmentation from
\texttt{OPTFM/node\_pretrain/main\_mip.py} lines 148--171).

\paragraph{Metrics.}
For each (model, transform, pair) cell we report mean cosine similarity
$\overline{\cos}(\Phi(G_A), \Phi(G_B))$, a 2{,}000-sample bootstrap 95\% CI,
and the fraction of samples where embeddings are bit-identical
(\texttt{exact\_frac}, computed via \texttt{torch.allclose} with
$\text{atol}\!=\!10^{-5}$).

\section{Main result: bit-identical embeddings on 1-WL pairs}
\label{sec:main}

\Cref{tab:main} reports the baseline, virtual-global-node, and
\rwpe{} columns of the main 6-model $\times$ 8-transform $\times$ 15-pair
grid (full grid in \texttt{results/main/20260410\_200438/report.txt}).
Every architecture, including the pretrained \sgformer{}+GCN and the full
random-init hierarchical \optfm{}, produces $\overline{\cos}\!=\!1.000000$
with bit-identical embeddings on every one of the 15 pairs. The
virtual-global-node augmentation (Lemma~\ref{lem:virtnode}) leaves this
unchanged. \rwpe{} (\Cref{sec:rwpe}) breaks the bit-identity for every
architecture, with $\texttt{exact\_frac}\!=\!0$.

\begin{table}[t]
\centering
\small
\caption{Main result. Each cell is the mean cosine similarity over the 15
1-WL-equivalent non-isomorphic pairs (Family I + Family II). Bracketed values
are 95\% bootstrap CIs; \texttt{exact \%} is the fraction of pairs where
$\Phi(G_A)$ and $\Phi(G_B)$ are bit-identical
($\|\cdot\|_\infty\!\le\!10^{-5}$). \textbf{Baseline}: theorem prediction holds
for every architecture. \textbf{VGN}: empirical confirmation of
\Cref{lem:virtnode}. \textbf{\rwpe{}}: positive control---every architecture
drops below 1.0; the bound \emph{is} falsifiable.}
\label{tab:main}
\begin{tabular}{@{}lrrrrr@{}}
\toprule
& \multicolumn{2}{c}{Baseline} & \multicolumn{1}{c}{VGN} & \multicolumn{2}{c}{\rwpe{} (4,6,8)} \\
\cmidrule(lr){2-3}\cmidrule(lr){4-4}\cmidrule(lr){5-6}
Model & $\overline{\cos}$ & exact \% & $\overline{\cos}$ & $\overline{\cos}$ & exact \% \\
\midrule
\sgformer{}+GCN (pretrained)        & 1.000000 & 100.0 & 1.000000 & 0.999486 & 0.0 \\
\sgformer{}+GCN (random)            & 1.000000 & 100.0 & 1.000000 & 0.998728 & 0.0 \\
\texttt{TransConvOnly} (random)     & 1.000000 & 100.0 & 1.000000 & 0.994253 & 0.0 \\
\texttt{GNNOnly} (random)           & 1.000000 & 100.0 & 1.000000 & 0.999994 & 0.0 \\
\texttt{SimpleGCN} (random)         & 1.000000 & 100.0 & 1.000000 & 0.994561 & 0.0 \\
\textbf{Hierarchical \optfm{} (random)} & \textbf{1.000000} & \textbf{100.0} & \textbf{1.000000} & \textbf{0.994022} & \textbf{0.0} \\
\bottomrule
\end{tabular}
\end{table}

\begin{figure}[t]
\centering
\includegraphics[width=\linewidth]{results/figures/fig2_main_result.pdf}
\caption{\textbf{1-WL boundedness is architecturally universal.}
(a) Per-cell distance from bit-identical, $1-\overline{\cos}(G_A,G_B)$,
on a logarithmic color scale; mean over 15 1-WL-equivalent non-isomorphic
pairs. Cells marked `=' hit bit-identical
($\overline{\cos}=1.000000$ to within $10^{-5}$). The Baseline and VGN
columns are `=' across all six architectures. \texttt{lp\_primal} is `='
for most models because variable features stay uniform on both graphs
(only the dual-feedback channel occasionally leaks signal). Only \rwpe{}
reliably produces a non-trivial gap. (b) Scale invariance. For every
$k\!\in\!\{2,\dots,30\}$ on Family I, all six architectures sit at
$\overline{\cos}=1.0$; the hierarchical \optfm{} under \rwpe{} (open
squares, dashed) is plotted as a falsifiability reference and drops below
1.0 at every scale.}
\label{fig:main}
\end{figure}

\subsection{Layer-wise verification}
\label{sec:layerwise}

The proof of \Cref{thm:main} proves the bound stage-by-stage. We instrument
the hierarchical \optfm{} and compute $\overline{\cos}(G_A, G_B)$ on the
pooled representation after each of its nine sub-layers
(\Cref{fig:layerwise}). The baseline curve sits at
$\overline{\cos}\!=\!1.0\pm 10^{-7}$ at every stage across 5 random
initializations, empirically verifying~\Cref{lem:linattn}--\Cref{lem:fuse}
individually. The VGN curve coincides at the same noise level, confirming
\Cref{lem:virtnode} layer-by-layer. The \rwpe{} curve drops below 1.0 at
the first sub-layer that sees the perturbed input (S1), propagates through
self-attention and cross-attention with genuine seed variance, and---notably
---recovers toward 1.0 at S6 because the GCN branch re-embeds from the raw
input via a parameter group whose random-init response is effectively blind
to the \rwpe{} channels. This recovery partially attenuates the signal after
fusion at S7 and pooling at S8, explaining why the end-to-end \rwpe{} drop
in~\Cref{tab:main} is modest in cosine.

\begin{figure}[t]
\centering
\includegraphics[width=0.85\linewidth]{results/figures/fig4_layerwise.pdf}
\caption{\textbf{The 1-WL bound holds at every sub-layer of the hierarchical
\optfm{}, not only end-to-end.} $\overline{\cos}(G_A,G_B)$ on the pooled
representation after each sub-layer, on the $k\!=\!10$ Family I pair.
Stages are grouped by the lemma they corroborate. Curves are the mean over
5 independent random initializations of the hierarchical \optfm{}.
\textbf{Baseline} (blue): $\overline{\cos}=1.0\pm 10^{-7}$ at every stage,
empirically verifying~\Cref{lem:linattn,lem:cross,lem:gcn,lem:fuse}.
\textbf{VGN} (green): coincides with baseline,
empirically confirming~\Cref{lem:virtnode}. \textbf{\rwpe{}} (red,
dashed; band = min--max over seeds): the signal enters at S1, propagates
through self/cross-attention, partially recovers at S6 because the GCN
branch re-embeds from the raw input.}
\label{fig:layerwise}
\end{figure}

\subsection{Training cannot overcome the bound}
\label{sec:plateau}

A natural reviewer concern is that the previous results use only random or
pretrained weights without optimization on the pair task. \Cref{thm:main}
rules out the possibility that training can find weights distinguishing a
1-WL-equivalent pair, but we verify this directly. We attach a 2-class
linear head to the mean-pooled hidden state of the hierarchical \optfm{}
and jointly train backbone+head as a binary classifier separating $G_A(k)$
(label 0) from $G_B(k)$ (label 1) on all 30 examples from the 15 pairs
(balanced 50/50). Cross-entropy loss, Adam at lr $10^{-3}$, full-batch,
300 epochs, 3 seeds (\texttt{scripts/finetune\_pair\_classifier.py}).

\paragraph{Theoretical floor.}
If \Cref{thm:main} holds, the pooled embedding is bit-identical for $G_A$
and $G_B$ at every step, the head outputs identical logits for both inputs
of every pair, and the loss is pinned at
$\mathcal L^\star = \tfrac{1}{2}(-\log p - \log(1-p))$ minimized at
$p=1/2$, giving $\mathcal L^\star = \ln 2 \approx 0.693147$.

\paragraph{Result (\Cref{fig:plateau}).}
Baseline final losses across seeds:
$0.693147,\,0.693147,\,0.693147$---matching $\ln 2$ to six decimal places.
Initial random-head losses were $0.6939, 0.6997, 0.7206$. Training is
moving the loss but only \emph{toward} the theoretical floor; it cannot
cross it. End-of-training pooled-embedding identity is preserved on every
one of the 15 pairs ($\texttt{frac\_pairs\_identical} = 1.000$). With \rwpe{}
on the same architecture and training procedure, the loss converges to
$2{\times}10^{-6}$, $2.7{\times}10^{-4}$, and $4.8{\times}10^{-4}$ across
the three seeds (5--7 decades below the baseline plateau), and the
embedding identity is broken on every pair
($\texttt{frac\_pairs\_identical} = 0$).

\begin{figure}[t]
\centering
\includegraphics[width=0.85\linewidth]{results/figures/fig5_training_plateau.pdf}
\caption{\textbf{Training cannot overcome the 1-WL bound.} Cross-entropy loss
vs.\ epoch for a 2-class linear-head classifier trained jointly with the
hierarchical \optfm{} backbone on the 15-pair 1-WL-equivalent population
(3 seeds per curve). \textbf{Baseline} (blue, solid): loss plateaus at
$\ln 2 \approx 0.693147$ to six decimals across all seeds for all 300 epochs;
end-of-training pooled embeddings remain bit-identical for every pair.
\textbf{\rwpe{}} (red, dashed): same architecture and training procedure
with \rwpe{}(4,6,8) inputs converges to $\sim\!3\!\times\!10^{-4}$ and breaks
the pooled-embedding identity on every pair.}
\label{fig:plateau}
\end{figure}

\section{Probe battery: structural and primal--dual targets}
\label{sec:probes}

Cosine similarity is a coarse metric. A finer probe asks: \emph{given frozen
\optfm{} embeddings, can a probe predict a structural or primal--dual
target?}

\paragraph{Datasets.}
$\mathcal D_{\text{rand}}$: 800 random bipartite MILPs at $(n_C, n_V)=(10,10)$,
sparsity $\approx\!0.25$, 50/50 equality/$\le$ constraints, random RHS and
objective; $\texttt{n\_components}\!\in\![1,5]$, $\texttt{lp\_value}\!\in\![-6.77, 4.52]$.
$\mathcal D_{\text{1wl}}$: 420 instances. For every $k\!\in\!\{2,\dots,30\}$ and
each of 15 random objective vectors, take both $C_{4k}$- and $k\!\cdot\!C_4$-
bipartite MILPs with that objective. All instances are mutually 1-WL
equivalent; $\texttt{n\_components}\!\in\![1,30]$ (\textbf{6$\times$ larger
variance} than $\mathcal D_{\text{rand}}$),
$\texttt{lp\_value}\!\in\![-29.16, 8.06]$.

\paragraph{Probes.}
Linear (\texttt{LinearProbe}) and 64-unit 2-hidden-layer MLP
(\texttt{MLPProbe}). 70/30 train/test split, Adam lr $10^{-2}$, 800 epochs,
5 seeds.

\begin{table}[t]
\centering
\small
\caption{Frozen-backbone probe scores. Targets: \texttt{n\_components}
(regression $R^2$), \texttt{lp\_value} (regression $R^2$), \texttt{girth\_le\_4}
and \texttt{feasible} (classification accuracy, majority baseline in parens).
The load-bearing row is \texttt{n\_components} (MLP): $R^2$ collapses from
${+}0.387$ on $\mathcal D_{\text{rand}}$ to ${-}0.011$ on $\mathcal
D_{\text{1wl}}$ \emph{despite the target having $6{\times}$ larger variance
on the latter}.}
\label{tab:probes}
\begin{tabular}{@{}lrr@{}}
\toprule
Target & $\mathcal D_{\text{rand}}$ (positive control) & $\mathcal D_{\text{1wl}}$ (target) \\
\midrule
\texttt{n\_components} linear $R^2$ & ${+}0.213\pm 0.060$ & ${-}0.012\pm 0.014$ \\
\texttt{n\_components} MLP $R^2$    & $\mathbf{{+}0.387 \pm 0.103}$ & $\mathbf{{-}0.011 \pm 0.014}$ \\
\texttt{lp\_value} linear $R^2$     & ${-}0.005\pm 0.022$ & ${-}0.015\pm 0.010$ \\
\texttt{lp\_value} MLP $R^2$        & ${-}0.000\pm 0.018$ & ${-}0.009\pm 0.011$ \\
\texttt{girth\_le\_4} linear acc    & 0.972 (maj 0.971)  & 0.449 (maj 0.551) \\
\texttt{girth\_le\_4} MLP acc       & 0.973 (maj 0.971)  & 0.449 (maj 0.551) \\
\texttt{feasible} linear acc        & 0.889 (maj 0.889)  & 1.000 (degenerate) \\
\texttt{feasible} MLP acc           & 0.889 (maj 0.889)  & 1.000 (degenerate) \\
\bottomrule
\end{tabular}
\end{table}

\begin{figure}[t]
\centering
\includegraphics[width=\linewidth]{results/figures/fig3_probe_battery.pdf}
\caption{\textbf{Frozen-backbone probe $R^2$ collapses on the 1-WL population
despite larger target variance.} (a) Linear and MLP probe scores for four
graph-level targets: $R^2$ for regression, accuracy-minus-majority for
classification. On random MILPs (positive control, blue), the MLP recovers
\texttt{n\_components} at $R^2\!=\!{+}0.40$. On the 1-WL-equivalent population
(red), the same probe collapses to $R^2\!=\!{-}0.01$. Error bars: $\pm 1$ std
over 5 seeds. \texttt{lp\_value} is at chance on both populations; \texttt{girth\_le\_4}
fits \emph{worse} than the majority classifier on $\mathcal D_{\text{1wl}}$,
consistent with embedding blindness; \texttt{feasible} is degenerate on
$\mathcal D_{\text{1wl}}$. (b) Variance sanity check: 1-WL targets have
$\ge$ as much variance as random ($6{\times}$ for \texttt{n\_components}),
ruling out a degenerate-target explanation.}
\label{fig:probes}
\end{figure}

\paragraph{Interpretation.}
The load-bearing cell is the second row of \Cref{tab:probes}. The collapse
from $R^2\!=\!{+}0.387$ to $R^2\!=\!{-}0.011$ happens even though
$\texttt{n\_components}$ varies six times more on $\mathcal D_{\text{1wl}}$
than on $\mathcal D_{\text{rand}}$---ruling out a degenerate-target
explanation. The $\texttt{lp\_value}$ rows give a complementary signal: on
random MILPs the embedding cannot predict LP value above chance, suggesting
the masked-edge pretraining objective does not implicitly induce LP-aware
representations. The $\texttt{girth\_le\_4}$ row is below the majority
baseline on $\mathcal D_{\text{1wl}}$, consistent with the embedding being
\emph{blind} to the distinguishing invariant.

\section{Falsifiability via \rwpe{}}
\label{sec:rwpe}

A purely null result invites the ``pathological pair'' objection: maybe
\emph{no} model can distinguish $C_{4k}$ from $k\cdot C_4$ regardless of
input. Random Walk Positional Encodings~\cite{dwivedi2022rwpe} are strictly
more expressive than 1-WL: at step $L\!=\!4$, $P^4[i,i] = 3/8$ on $C_{4k}$
($k\!\ge\!2$) but $P^4[i,i] = 1/2$ on $k\cdot C_4$. We inject return
probabilities at $L\!\in\!\{4,6,8\}$ into the three zero-valued Ecole
variable feature slots, plus $L\!=\!4$ into the constraint feature
(\texttt{scripts/improvements.make\_rwpe\_transform}).

The \rwpe{} columns of \Cref{tab:main} show every architecture drops below
$\overline{\cos}\!=\!1.0$, with $\texttt{exact\_frac}$ falling from 100\%
to 0\% on every architecture and both families. The drops are largest for
the simpler GCN-style architectures and the hierarchical \optfm{}; the
randomly-initialized GNN-only ablation barely moves because its random
weight matrices happen to attenuate the \rwpe{} input signal.

The composition of this falsifiability test with the layer-wise probe
(\Cref{fig:layerwise}) and the supervised fine-tune (\Cref{fig:plateau})
makes the strongest possible empirical case. The architectural bound holds
end-to-end, at every sub-layer, and survives 300 epochs of optimization;
yet a strictly more-expressive input transform breaks all three
simultaneously.

\section{Discussion}
\label{sec:discussion}

\paragraph{Relation to prior work.}
Chen et al.~\cite{chen2023expressive} proved the 1-WL bound for MPNNs on
bipartite MILPs and exhibited a pair where feasibility differs across the
1-WL class. We extend their argument to an architecture explicitly motivated
as going \emph{beyond} GNN-style local interaction: \optfm{}'s hierarchical
multi-view cross-attention. Our two new lemmas
(\Cref{lem:linattn,lem:virtnode}) cover the architectural novelty---global
linear attention with edge-weighted cross-attention, and the
virtual-global-node pretraining augmentation---which were not analyzed in
prior work. Joshi et al.~\cite{joshi2024} ported the 1-WL argument to
geometric GNNs; ours is the corresponding port to combinatorial
optimization foundation models.

\paragraph{What we have not established.}
The theorem is an expressiveness statement, not a utility statement. It
does not contradict \optfm{}'s reported gains on CA, MIS, MVC, and SC
benchmarks~\cite{yuan2025optfm}, because those tasks admit 1-WL solutions on
generic inputs---the bound is not the binding constraint there. We have not
retrained \optfm{} with \rwpe{} features (the pretraining dataset is not
fully public and requires multi-A100 compute) nor extended the formal
analysis to \optfm{}'s graph-level contrastive pretraining stage, which
operates on subgraph tokens orthogonal to the node-level expressiveness
question studied here. Finally, our pair families are minimum viable
1-WL tests; a journal version would extend to Cai--F\"urer--Immerman
constructions over a continuous distribution of 1-WL classes.

\paragraph{Implication for foundation models in optimization.}
\optfm{}'s headline architectural claim---that multi-view linear attention
plus edge-weighted cross-attention enhances representational power
beyond local GNN interaction---is not borne out at the level of expressiveness.
Any practical gains attributable to \optfm{}'s architecture must therefore
arise from optimization dynamics, inductive biases for tasks within the 1-WL
class, or downstream transfer effects, \emph{not} from extra discriminative
power on bipartite MILPs. Future foundation-model designs targeting
combinatorial optimization should include positional encodings or random
features at the input level if expressiveness beyond 1-WL is required.

\begin{thebibliography}{99}\small

\bibitem{gasse2019exact}
M.~Gasse, D.~Ch\'etelat, N.~Ferroni, L.~Charlin, and A.~Lodi.
Exact combinatorial optimization with graph convolutional neural networks.
In \emph{NeurIPS}, 2019.

\bibitem{chen2023expressive}
Z.~Chen, J.~Liu, X.~Wang, J.~Lu, and W.~Yin.
On representing mixed-integer linear programs by graph neural networks.
In \emph{ICLR}, 2023.

\bibitem{yuan2025optfm}
H.~Yuan, W.~Ouyang, C.~Zhang, C.~Li, and Y.~Sun.
{OPTFM}: A scalable multi-view graph transformer for hierarchical
pre-training in combinatorial optimization.
In \emph{NeurIPS}, 2025.

\bibitem{xu2019powerful}
K.~Xu, W.~Hu, J.~Leskovec, and S.~Jegelka.
How powerful are graph neural networks?
In \emph{ICLR}, 2019.

\bibitem{morris2019wlgo}
C.~Morris, M.~Ritzert, M.~Fey, W.~L.~Hamilton, J.~E.~Lenssen, G.~Rattan, and M.~Grohe.
Weisfeiler and {L}eman go neural: Higher-order graph neural networks.
In \emph{AAAI}, 2019.

\bibitem{wu2023sgformer}
Q.~Wu, C.~Yang, W.~Zhao, Y.~He, D.~Wipf, and J.~Yan.
{SGFormer}: Simplifying and empowering transformers for large-graph
representations.
In \emph{NeurIPS}, 2023.

\bibitem{dwivedi2022rwpe}
V.~P.~Dwivedi, C.~K.~Joshi, A.~T.~Luu, T.~Laurent, Y.~Bengio, and X.~Bresson.
Graph neural networks with learnable structural and positional representations.
In \emph{ICLR}, 2022.

\bibitem{abboud2020rni}
R.~Abboud, \.{I}.~\.{I}.~Ceylan, M.~Grohe, and T.~Lukasiewicz.
The surprising power of graph neural networks with random node initialization.
In \emph{IJCAI}, 2020.

\bibitem{sato2021random}
R.~Sato, M.~Yamada, and H.~Kashima.
Random features strengthen graph neural networks.
In \emph{SDM}, 2021.

\bibitem{cai1992lower}
J.-Y. Cai, M.~F\"urer, and N.~Immerman.
An optimal lower bound on the number of variables for graph identification.
\emph{Combinatorica}, 12(4):389--410, 1992.

\bibitem{joshi2024}
C.~K.~Joshi, C.~Bodnar, S.~V.~Mathis, T.~Cohen, and P.~Li\'o.
On the expressive power of geometric graph neural networks.
In \emph{ICML}, 2024.

\bibitem{khalil2017learning}
E.~Khalil, P.~L.~Bodic, L.~Song, G.~Nemhauser, and B.~Dilkina.
Learning to branch in mixed integer programming.
In \emph{AAAI}, 2017.

\bibitem{gupta2020hybrid}
P.~Gupta, M.~Gasse, E.~Khalil, P.~Mudigonda, A.~Lodi, and Y.~Bengio.
Hybrid models for learning to branch.
In \emph{NeurIPS}, 2020.

\end{thebibliography}

\appendix

\section{The deprecated ``canonical 1-WL pair''}
\label{app:broken-canonical}

The ``canonical 1-WL pair'' on 6 variables and 5 constraints, adapted
from~\cite{chen2023expressive} in earlier probing code (\texttt{data/milp\_pairs.construct\_canonical\_pair}),
turns out to be \emph{graph-isomorphic} as a feature-labeled bipartite
graph, not 1-WL equivalent non-isomorphic. We document the diagnosis
(\texttt{scripts/sanity\_check\_pairs.py}):
\begin{itemize}[leftmargin=18pt, topsep=0pt, itemsep=1pt]
    \item MILP-A is feasible at $(1,0,0,1,0,1)$.
    \item MILP-B---which the original docstring claims is infeasible---is
    feasible at $(1,0,1,0,0,1)$ as found by \texttt{scipy.milp}. \emph{Both}
    MILPs are feasible.
    \item Brute-force permutation search produces an explicit isomorphism:
    constraint permutation $\pi=(0,1,2,4,3)$ and variable permutation
    $\sigma=(1,0,3,2,5,4)$ give $A[\pi,\sigma]=A$ and $b_A[\pi]=b_B$.
\end{itemize}
Any permutation-equivariant model trivially maps isomorphic inputs to
identical embeddings. The $\overline{\cos}\!=\!1.0$ result on this pair
from any GNN is therefore expected and uninformative as evidence of 1-WL
boundedness. This finding is the reason for constructing the corrected
pair families of \Cref{sec:setup}.

\section{Reproducibility}
\label{app:repro}

\paragraph{Environment.}
Windows 11 Pro 10.0.26200, bash via Git Bash; Python 3.14.2 in a venv;
\texttt{torch 2.11.0+cpu}, \texttt{numpy 2.4.4}, \texttt{scipy 1.17.1},
\texttt{matplotlib 3.10.8}, \texttt{networkx 3.6.1}. Working directory:
\texttt{d:/GitHub/optfm\_probing/}. The pretrained \optfm{} checkpoint is
read from \texttt{D:/GitHub/OPTFM/node\_pretrain/Models\_SCIP/model\_params\_epoch\_31.pth};
update \texttt{CKPT} constants if cloning OPTFM elsewhere. Set
\texttt{PYTHONUTF8=1} to handle the \checkmark and $\times$ characters
emitted by scripts.

\paragraph{Commands.}
\begin{lstlisting}[language=bash]
# 1. verify checkpoint loads with strict=True
PYTHONUTF8=1 python scripts/verify_weights.py

# 2. diagnose the deprecated canonical pair
PYTHONUTF8=1 python scripts/sanity_check_pairs.py

# 3. main 6 x 8 grid on Family I + Family II
PYTHONUTF8=1 python scripts/run_main_experiment.py \
    --k_values 2 3 4 5 6 7 8 9 10 12 15 20 25 30 \
    --rnf_samples 30 --family both

# 4. structural / primal-dual probe battery
PYTHONUTF8=1 python scripts/probe_primal_dual.py \
    --n_random 800 --k_values 2 3 4 5 6 7 8 9 10 12 15 20 25 30 \
    --n_objectives 15 --n_seeds 5

# 5. supervised fine-tune (Section 5.2 / Figure 5)
PYTHONUTF8=1 python scripts/finetune_pair_classifier.py \
    --epochs 300 --lr 1e-3 --seeds 0 1 2

# 6. regenerate figures
PYTHONUTF8=1 python visualization/generate_all.py
\end{lstlisting}

\paragraph{Outputs.}
\texttt{results/main/<ts>/\{results.json, summary.csv, report.txt\}}
(main grid);
\texttt{results/probes/<ts>/\{results.json, report.txt\}} (probe battery);
\texttt{results/finetune/<ts>/results.json} (training plateau);
\texttt{results/figures/fig\{1..5\}.\{pdf,png\}} (\Cref{fig:families,fig:main,fig:probes,fig:layerwise,fig:plateau}).

\paragraph{Repository layout.}
\begin{lstlisting}[basicstyle=\ttfamily\scriptsize]
optfm_probing/
|-- paper/{proof_sketch.md, FINAL_REPORT.md}
|-- models/{sgformer_mip.py, optfm_hierarchical.py}
|-- data/{milp_pairs.py [DEPRECATED], milp_pairs_v2.py}
|-- scripts/
|   |-- verify_weights.py            checkpoint integrity
|   |-- sanity_check_pairs.py        diagnoses the broken canonical pair
|   |-- run_main_experiment.py       6 x 8 main grid (Section 5)
|   |-- probe_primal_dual.py         probe battery (Section 6)
|   |-- finetune_pair_classifier.py  supervised plateau (Section 5.2)
|   |-- improvements.py              input transforms (RNF/LP/RWPE/VGN)
|   |-- _tee_log.py                  stdout tee
|-- visualization/
|   |-- neurips_style.py             matplotlib rcParams
|   |-- fig{1..5}_*.py               per-figure scripts
|   |-- generate_all.py              regenerates results/figures/
|-- results/{main, probes, finetune, figures, logs}/
|-- Project_Proposal.tex             this document
\end{lstlisting}

\end{document}
